\pagebreak
    \hspace{0pt}
    \vfill
    \addcontentsline{toc}{section}{Abstract}
    \section*{Abstract}
    
    In drug-resistant focal epilepsy, surgical treatment represents the best
    chance to reduce seizure frequency, and accurate localization of the
    Seizure Onset Zone (SOZ) is critical for pre-surgery planning.
    However, the high inter-patient heterogeneity of stereo-electroencephalography
    (sEEG) recordings makes the development of generalizable, patient-independent
    methods particularly challenging.

    This study proposes an automated pipeline for the localization of the SOZ
    across multiple patients using a unified parameter set.
    A total of 23 patients was included in the dataset, including resting-state
    and seizure sEEG recordings.
    22 features are extracted from data:
    signal features characterizing the time domain, spectral features capturing
    information from the frequency domain and network features modelling the functional
    relationships between brain regions.

    A Variational Autoencoder (VAE) produces a latent representation of the
    resting-state features, while a Gaussian Mixture Model (GMM) fits the
    encoded distribution and scores the likelihood of seizure features.
    A classifier decides which brain regions compose the SOZ based on the
    negative log-likelihood of seizure features, which is used as an indicator
    for anomaly detection.
    Hyperparameter optimization is performed including a preliminary validation
    phase using Leave-One-Out Cross-Validation (LOOCV).
    The final model is then trained on the full training set and evaluated
    on an independent test set of unseen patients.

    The proposed model achieved promising results during cross-validation,
    however a significant performance drop was observed on unseen test patients,
    reflecting the intrinsic heterogeneity of sEEG recordings across subjects.
    Results highlight that inter-patient variability remains the primary
    challenge in SOZ localization, however potential strategies are proposed
    to direct future work towards patient-adaptive approaches.
    \vfill
    \hspace{0pt}
\pagebreak

